%%%%%%%%%%%%%%%%%%%%%%%%%%%%%%%%%%%%%%%%%%%%%%%%%%%%%%%%%%%%
% 10.1 OPTIMIZATION FUNDAMENTALS
% The Composition of Advanced Mathematics
%%%%%%%%%%%%%%%%%%%%%%%%%%%%%%%%%%%%%%%%%%%%%%%%%%%%%%%%%%%%

\section{Optimization Fundamentals}
\label{sec:optimization-fundamentals}

\prerequisite{
Functions, algebra, single-variable and multivariable calculus,
gradients, Hessian matrices, linear algebra, systems of equations,
inequalities, and basic mathematical modeling.
}

\learningobjectives{
By the end of this section, the reader should be able to:
\begin{itemize}
    \item define an optimization problem;
    \item identify decision variables, objective functions, and
          constraints;
    \item distinguish minimization and maximization problems;
    \item define feasible sets and feasible points;
    \item distinguish local and global optima;
    \item define strict and non-strict optimality;
    \item recognize unconstrained and constrained optimization problems;
    \item formulate optimization models from verbal descriptions;
    \item understand the role of parameters and data;
    \item use first-order necessary conditions;
    \item use second-order derivative tests;
    \item understand positive definite and semidefinite Hessians;
    \item identify stationary and critical points;
    \item understand coercivity and existence of minimizers conceptually;
    \item recognize convex sets and convex functions;
    \item understand why convexity simplifies optimization;
    \item recognize affine and linear constraints;
    \item define active and inactive inequality constraints;
    \item understand equality-constrained optimization;
    \item derive Lagrange multiplier conditions;
    \item interpret Lagrange multipliers;
    \item understand constraint qualifications conceptually;
    \item state basic Karush--Kuhn--Tucker conditions;
    \item understand complementary slackness;
    \item distinguish necessary and sufficient optimality conditions;
    \item understand sensitivity and shadow-price interpretations;
    \item recognize optimization in vector and matrix notation;
    \item understand regularization conceptually;
    \item recognize least-squares optimization;
    \item understand gradient descent;
    \item recognize Newton's method for optimization;
    \item distinguish line-search and trust-region ideas conceptually;
    \item understand conditioning and scaling;
    \item understand deterministic versus stochastic optimization
          conceptually;
    \item connect optimization with operations research, machine learning,
          statistics, control, economics, and numerical mathematics.
\end{itemize}
}

%%%%%%%%%%%%%%%%%%%%%%%%%%%%%%%%%%%%%%%%%%%%%%%%%%%%%%%%%%%%
\subsection{What Is Optimization?}
\label{subsec:what-is-optimization}
%%%%%%%%%%%%%%%%%%%%%%%%%%%%%%%%%%%%%%%%%%%%%%%%%%%%%%%%%%%%

Optimization is the mathematical study of selecting the best feasible
choice according to a specified objective.

A general minimization problem has form

\[
\boxed{
\min_{\mathbf{x}}
f(\mathbf{x}),
}
\]

where

\[
\mathbf{x}
\]

contains the decision variables and

\[
f
\]

is the objective function.

%%%%%%%%%%%%%%%%%%%%%%%%%%%%%%%%%%%%%%%%%%%%%%%%%%%%%%%%%%%%
\subsection{Decision Variables}
\label{subsec:decision-variables}
%%%%%%%%%%%%%%%%%%%%%%%%%%%%%%%%%%%%%%%%%%%%%%%%%%%%%%%%%%%%

\begin{definitionbox}{Decision Variables}
Decision variables are quantities whose values are chosen by the
optimization process.
\end{definitionbox}

For example,

\[
\mathbf{x}
=
\begin{pmatrix}
x_1\\
x_2\\
\vdots\\
x_n
\end{pmatrix}
\in\mathbb{R}^n
\]

may represent:

\[
\boxed{
\text{production quantities},
}
\]

\[
\boxed{
\text{resource allocations},
}
\]

\[
\boxed{
\text{portfolio weights},
}
\]

or

\[
\boxed{
\text{model parameters}.
}
\]

%%%%%%%%%%%%%%%%%%%%%%%%%%%%%%%%%%%%%%%%%%%%%%%%%%%%%%%%%%%%
\subsection{Objective Function}
\label{subsec:objective-function}
%%%%%%%%%%%%%%%%%%%%%%%%%%%%%%%%%%%%%%%%%%%%%%%%%%%%%%%%%%%%

\begin{definitionbox}{Objective Function}
The objective function assigns a numerical value to each possible
decision.
\end{definitionbox}

It may represent:

\[
\boxed{
\text{cost},
}
\]

\[
\boxed{
\text{profit},
}
\]

\[
\boxed{
\text{error},
}
\]

\[
\boxed{
\text{distance},
}
\]

or

\[
\boxed{
\text{risk}.
}
\]

%%%%%%%%%%%%%%%%%%%%%%%%%%%%%%%%%%%%%%%%%%%%%%%%%%%%%%%%%%%%
\subsection{Minimization and Maximization}
\label{subsec:minimization-maximization}
%%%%%%%%%%%%%%%%%%%%%%%%%%%%%%%%%%%%%%%%%%%%%%%%%%%%%%%%%%%%

A minimization problem is

\[
\boxed{
\min_{\mathbf{x}}f(\mathbf{x}).
}
\]

A maximization problem is

\[
\boxed{
\max_{\mathbf{x}}f(\mathbf{x}).
}
\]

Every maximization problem can be converted into minimization:

\[
\boxed{
\max f(\mathbf{x})
\quad\Longleftrightarrow\quad
\min -f(\mathbf{x}).
}
\]

%%%%%%%%%%%%%%%%%%%%%%%%%%%%%%%%%%%%%%%%%%%%%%%%%%%%%%%%%%%%
\subsection{Constraints}
\label{subsec:optimization-constraints}
%%%%%%%%%%%%%%%%%%%%%%%%%%%%%%%%%%%%%%%%%%%%%%%%%%%%%%%%%%%%

Constraints restrict which decisions are allowed.

Equality constraints have form

\[
\boxed{
h_j(\mathbf{x})=0,
\qquad
j=1,\ldots,m.
}
\]

Inequality constraints are commonly written

\[
\boxed{
g_i(\mathbf{x})\leq0,
\qquad
i=1,\ldots,p.
}
\]

A general constrained problem is

\[
\boxed{
\begin{aligned}
\min_{\mathbf{x}}
\quad&
f(\mathbf{x})
\\
\text{subject to}
\quad&
g_i(\mathbf{x})\leq0,
\qquad
i=1,\ldots,p,
\\
&
h_j(\mathbf{x})=0,
\qquad
j=1,\ldots,m.
\end{aligned}
}
\]

%%%%%%%%%%%%%%%%%%%%%%%%%%%%%%%%%%%%%%%%%%%%%%%%%%%%%%%%%%%%
\subsection{Feasible Set}
\label{subsec:feasible-set}
%%%%%%%%%%%%%%%%%%%%%%%%%%%%%%%%%%%%%%%%%%%%%%%%%%%%%%%%%%%%

\begin{definitionbox}{Feasible Set}
The feasible set is the set of all decision vectors satisfying every
constraint.
\end{definitionbox}

If

\[
g_i(\mathbf{x})\leq0
\]

and

\[
h_j(\mathbf{x})=0,
\]

then

\[
\boxed{
\mathcal{F}
=
\left\{
\mathbf{x}:
g_i(\mathbf{x})\leq0,
\;
h_j(\mathbf{x})=0
\right\}.
}
\]

%%%%%%%%%%%%%%%%%%%%%%%%%%%%%%%%%%%%%%%%%%%%%%%%%%%%%%%%%%%%
\subsection{Feasible Points}
\label{subsec:feasible-points}
%%%%%%%%%%%%%%%%%%%%%%%%%%%%%%%%%%%%%%%%%%%%%%%%%%%%%%%%%%%%

A point

\[
\mathbf{x}\in\mathcal{F}
\]

is called feasible.

A point violating at least one constraint is infeasible.

Optimization therefore asks:

\[
\boxed{
\text{Which feasible point has the best objective value?}
}
\]

%%%%%%%%%%%%%%%%%%%%%%%%%%%%%%%%%%%%%%%%%%%%%%%%%%%%%%%%%%%%
\subsection{Optimal Solution}
\label{subsec:optimal-solution}
%%%%%%%%%%%%%%%%%%%%%%%%%%%%%%%%%%%%%%%%%%%%%%%%%%%%%%%%%%%%

A feasible point \(\mathbf{x}^*\) is a global minimizer if

\[
\boxed{
f(\mathbf{x}^*)
\leq
f(\mathbf{x})
}
\]

for every

\[
\mathbf{x}\in\mathcal{F}.
\]

The optimal objective value is

\[
\boxed{
f^*
=
f(\mathbf{x}^*).
}
\]

%%%%%%%%%%%%%%%%%%%%%%%%%%%%%%%%%%%%%%%%%%%%%%%%%%%%%%%%%%%%
\subsection{Local Minimizer}
\label{subsec:local-minimizer}
%%%%%%%%%%%%%%%%%%%%%%%%%%%%%%%%%%%%%%%%%%%%%%%%%%%%%%%%%%%%

A feasible point \(\mathbf{x}^*\) is a local minimizer if there exists a
neighborhood around \(\mathbf{x}^*\) such that

\[
\boxed{
f(\mathbf{x}^*)
\leq
f(\mathbf{x})
}
\]

for all feasible \(\mathbf{x}\) sufficiently close to
\(\mathbf{x}^*\).

A local solution need not be globally optimal.

%%%%%%%%%%%%%%%%%%%%%%%%%%%%%%%%%%%%%%%%%%%%%%%%%%%%%%%%%%%%
\subsection{Strict Local Minimizer}
\label{subsec:strict-local-minimizer}
%%%%%%%%%%%%%%%%%%%%%%%%%%%%%%%%%%%%%%%%%%%%%%%%%%%%%%%%%%%%

A strict local minimizer satisfies

\[
\boxed{
f(\mathbf{x}^*)
<
f(\mathbf{x})
}
\]

for every nearby feasible point

\[
\mathbf{x}\neq\mathbf{x}^*.
\]

Strictness concerns the objective values of neighboring feasible points.

%%%%%%%%%%%%%%%%%%%%%%%%%%%%%%%%%%%%%%%%%%%%%%%%%%%%%%%%%%%%
\subsection{Global Versus Local Optimization}
\label{subsec:global-versus-local-optimization}
%%%%%%%%%%%%%%%%%%%%%%%%%%%%%%%%%%%%%%%%%%%%%%%%%%%%%%%%%%%%

Local optimization asks for the best point within a neighborhood.

Global optimization asks for the best point in the entire feasible set.

For nonconvex problems,

\[
\boxed{
\text{local optimum}
\not\Rightarrow
\text{global optimum}.
}
\]

This distinction is one of the central difficulties of optimization.

%%%%%%%%%%%%%%%%%%%%%%%%%%%%%%%%%%%%%%%%%%%%%%%%%%%%%%%%%%%%
\subsection{Unconstrained Optimization}
\label{subsec:unconstrained-optimization}
%%%%%%%%%%%%%%%%%%%%%%%%%%%%%%%%%%%%%%%%%%%%%%%%%%%%%%%%%%%%

An unconstrained problem has form

\[
\boxed{
\min_{\mathbf{x}\in\mathbb{R}^n}
f(\mathbf{x}).
}
\]

The only restriction is that the decision vector belongs to its domain.

Differential calculus provides basic local optimality conditions.

%%%%%%%%%%%%%%%%%%%%%%%%%%%%%%%%%%%%%%%%%%%%%%%%%%%%%%%%%%%%
\subsection{Constrained Optimization}
\label{subsec:constrained-optimization}
%%%%%%%%%%%%%%%%%%%%%%%%%%%%%%%%%%%%%%%%%%%%%%%%%%%%%%%%%%%%

A constrained problem restricts the decision vector:

\[
\boxed{
\min_{\mathbf{x}\in\mathcal{F}}
f(\mathbf{x}).
}
\]

The optimum may occur:

\[
\boxed{
\text{inside the feasible region},
}
\]

or

\[
\boxed{
\text{on its boundary}.
}
\]

Constraint geometry therefore becomes important.

%%%%%%%%%%%%%%%%%%%%%%%%%%%%%%%%%%%%%%%%%%%%%%%%%%%%%%%%%%%%
\subsection{Optimization Model}
\label{subsec:optimization-model}
%%%%%%%%%%%%%%%%%%%%%%%%%%%%%%%%%%%%%%%%%%%%%%%%%%%%%%%%%%%%

A mathematical optimization model contains three central components:

\[
\boxed{
\text{decision variables},
}
\]

\[
\boxed{
\text{objective function},
}
\]

and

\[
\boxed{
\text{constraints}.
}
\]

A fourth important component is the model data or parameters.

%%%%%%%%%%%%%%%%%%%%%%%%%%%%%%%%%%%%%%%%%%%%%%%%%%%%%%%%%%%%
\subsection{Worked Example: Basic Model Formulation}
\label{subsec:worked-basic-optimization-model}
%%%%%%%%%%%%%%%%%%%%%%%%%%%%%%%%%%%%%%%%%%%%%%%%%%%%%%%%%%%%

\begin{workedexamplebox}{Production Planning}

Suppose a company produces products \(A\) and \(B\).

Let

\[
x
=
\text{units of product }A,
\]

and

\[
y
=
\text{units of product }B.
\]

Suppose profit is

\[
5
\]

per unit of \(A\) and

\[
8
\]

per unit of \(B\).

The objective is

\[
\boxed{
\max
\;
5x+8y.
}
\]

Suppose resource limits are

\[
2x+y\leq100,
\]

and

\[
x+3y\leq120.
\]

Nonnegativity requires

\[
x\geq0,
\qquad
y\geq0.
\]

Therefore the optimization model is

\begin{answerbox}
\[
\boxed{
\begin{aligned}
\max_{x,y}
\quad&
5x+8y
\\
\text{subject to}
\quad&
2x+y\leq100,
\\
&
x+3y\leq120,
\\
&
x\geq0,
\quad
y\geq0.
\end{aligned}
}
\]
\end{answerbox}

\end{workedexamplebox}

%%%%%%%%%%%%%%%%%%%%%%%%%%%%%%%%%%%%%%%%%%%%%%%%%%%%%%%%%%%%
\subsection{Parameters and Data}
\label{subsec:optimization-parameters-data}
%%%%%%%%%%%%%%%%%%%%%%%%%%%%%%%%%%%%%%%%%%%%%%%%%%%%%%%%%%%%

Parameters are known quantities entering the model.

For example,

\[
\boxed{
c_i
=
\text{unit costs},
}
\]

\[
\boxed{
a_{ij}
=
\text{resource requirements},
}
\]

and

\[
\boxed{
b_j
=
\text{available capacities}.
}
\]

The optimizer chooses the variables, not the parameters.

%%%%%%%%%%%%%%%%%%%%%%%%%%%%%%%%%%%%%%%%%%%%%%%%%%%%%%%%%%%%
\subsection{First-Order Necessary Condition}
\label{subsec:first-order-necessary-condition}
%%%%%%%%%%%%%%%%%%%%%%%%%%%%%%%%%%%%%%%%%%%%%%%%%%%%%%%%%%%%

Suppose \(f:\mathbb{R}^n\to\mathbb{R}\) is differentiable and
\(\mathbf{x}^*\) is an unconstrained local minimizer lying in the
interior of the domain.

Then

\[
\boxed{
\nabla f(\mathbf{x}^*)
=
\mathbf{0}.
}
\]

This is a necessary condition, not generally sufficient.

%%%%%%%%%%%%%%%%%%%%%%%%%%%%%%%%%%%%%%%%%%%%%%%%%%%%%%%%%%%%
\subsection{Stationary Points}
\label{subsec:stationary-points-optimization}
%%%%%%%%%%%%%%%%%%%%%%%%%%%%%%%%%%%%%%%%%%%%%%%%%%%%%%%%%%%%

A stationary point satisfies

\[
\boxed{
\nabla f(\mathbf{x})=\mathbf{0}.
}
\]

A stationary point may be:

\[
\boxed{
\text{a local minimum},
}
\]

\[
\boxed{
\text{a local maximum},
}
\]

or

\[
\boxed{
\text{a saddle point}.
}
\]

%%%%%%%%%%%%%%%%%%%%%%%%%%%%%%%%%%%%%%%%%%%%%%%%%%%%%%%%%%%%
\subsection{Gradient}
\label{subsec:optimization-gradient}
%%%%%%%%%%%%%%%%%%%%%%%%%%%%%%%%%%%%%%%%%%%%%%%%%%%%%%%%%%%%

For

\[
f(x_1,\ldots,x_n),
\]

the gradient is

\[
\boxed{
\nabla f
=
\begin{pmatrix}
\dfrac{\partial f}{\partial x_1}
\\[0.5em]
\vdots
\\[0.5em]
\dfrac{\partial f}{\partial x_n}
\end{pmatrix}.
}
\]

The gradient points in the direction of steepest local increase under the
standard Euclidean geometry.

Therefore

\[
\boxed{
-\nabla f
}
\]

is the direction of steepest local decrease.

%%%%%%%%%%%%%%%%%%%%%%%%%%%%%%%%%%%%%%%%%%%%%%%%%%%%%%%%%%%%
\subsection{Directional Derivative}
\label{subsec:optimization-directional-derivative}
%%%%%%%%%%%%%%%%%%%%%%%%%%%%%%%%%%%%%%%%%%%%%%%%%%%%%%%%%%%%

For a direction \(\mathbf{d}\), the directional derivative is

\[
\boxed{
D_{\mathbf{d}}f(\mathbf{x})
=
\nabla f(\mathbf{x})^T\mathbf{d}.
}
\]

If

\[
\nabla f^T\mathbf{d}<0,
\]

then \(\mathbf{d}\) is a descent direction.

%%%%%%%%%%%%%%%%%%%%%%%%%%%%%%%%%%%%%%%%%%%%%%%%%%%%%%%%%%%%
\subsection{Worked Example: Stationary Point}
\label{subsec:worked-stationary-point}
%%%%%%%%%%%%%%%%%%%%%%%%%%%%%%%%%%%%%%%%%%%%%%%%%%%%%%%%%%%%

\begin{workedexamplebox}{Quadratic Objective}

Find the stationary point of

\[
f(x,y)
=
x^2+2y^2-4x+8y.
\]

Compute

\[
f_x
=
2x-4,
\]

and

\[
f_y
=
4y+8.
\]

Set both equal to zero:

\[
2x-4=0,
\]

\[
4y+8=0.
\]

Thus,

\[
x=2,
\qquad
y=-2.
\]

Therefore,

\begin{answerbox}
\[
\boxed{
\mathbf{x}^*
=
(2,-2).
}
\]
\end{answerbox}

\end{workedexamplebox}

%%%%%%%%%%%%%%%%%%%%%%%%%%%%%%%%%%%%%%%%%%%%%%%%%%%%%%%%%%%%
\subsection{Hessian Matrix}
\label{subsec:hessian-optimization}
%%%%%%%%%%%%%%%%%%%%%%%%%%%%%%%%%%%%%%%%%%%%%%%%%%%%%%%%%%%%

The Hessian matrix is

\[
\boxed{
\nabla^2f(\mathbf{x})
=
\begin{pmatrix}
\dfrac{\partial^2f}{\partial x_1^2}
&
\cdots
&
\dfrac{\partial^2f}{\partial x_1\partial x_n}
\\
\vdots
&
\ddots
&
\vdots
\\
\dfrac{\partial^2f}{\partial x_n\partial x_1}
&
\cdots
&
\dfrac{\partial^2f}{\partial x_n^2}
\end{pmatrix}.
}
\]

For sufficiently smooth \(f\), the Hessian is symmetric.

%%%%%%%%%%%%%%%%%%%%%%%%%%%%%%%%%%%%%%%%%%%%%%%%%%%%%%%%%%%%
\subsection{Positive Definite Matrices}
\label{subsec:positive-definite-optimization}
%%%%%%%%%%%%%%%%%%%%%%%%%%%%%%%%%%%%%%%%%%%%%%%%%%%%%%%%%%%%

A symmetric matrix \(H\) is positive definite if

\[
\boxed{
\mathbf{d}^TH\mathbf{d}>0
}
\]

for every nonzero vector \(\mathbf{d}\).

It is positive semidefinite if

\[
\boxed{
\mathbf{d}^TH\mathbf{d}\geq0.
}
\]

%%%%%%%%%%%%%%%%%%%%%%%%%%%%%%%%%%%%%%%%%%%%%%%%%%%%%%%%%%%%
\subsection{Second-Order Sufficient Condition}
\label{subsec:second-order-sufficient-condition}
%%%%%%%%%%%%%%%%%%%%%%%%%%%%%%%%%%%%%%%%%%%%%%%%%%%%%%%%%%%%

Suppose

\[
\nabla f(\mathbf{x}^*)=0
\]

and

\[
\boxed{
\nabla^2f(\mathbf{x}^*)
\succ0.
}
\]

Then \(\mathbf{x}^*\) is a strict local minimizer.

Similarly, if

\[
\boxed{
\nabla^2f(\mathbf{x}^*)
\prec0,
}
\]

then \(\mathbf{x}^*\) is a strict local maximizer.

%%%%%%%%%%%%%%%%%%%%%%%%%%%%%%%%%%%%%%%%%%%%%%%%%%%%%%%%%%%%
\subsection{Indefinite Hessian}
\label{subsec:indefinite-hessian}
%%%%%%%%%%%%%%%%%%%%%%%%%%%%%%%%%%%%%%%%%%%%%%%%%%%%%%%%%%%%

If the Hessian is indefinite, there exist directions
\(\mathbf{d}_1,\mathbf{d}_2\) such that

\[
\mathbf{d}_1^TH\mathbf{d}_1>0,
\]

while

\[
\mathbf{d}_2^TH\mathbf{d}_2<0.
\]

At a stationary point, this indicates saddle-type behavior.

%%%%%%%%%%%%%%%%%%%%%%%%%%%%%%%%%%%%%%%%%%%%%%%%%%%%%%%%%%%%
\subsection{Worked Example: Hessian Test}
\label{subsec:worked-hessian-test}
%%%%%%%%%%%%%%%%%%%%%%%%%%%%%%%%%%%%%%%%%%%%%%%%%%%%%%%%%%%%

\begin{workedexamplebox}{Classifying a Stationary Point}

For

\[
f(x,y)
=
x^2+2y^2-4x+8y,
\]

the Hessian is

\[
\nabla^2f
=
\begin{pmatrix}
2&0\\
0&4
\end{pmatrix}.
\]

Its eigenvalues are

\[
2,
\qquad
4.
\]

Both are positive.

Therefore the Hessian is positive definite.

Hence,

\begin{answerbox}
\[
\boxed{
(2,-2)
\text{ is a strict local minimizer}.
}
\]
\end{answerbox}

Because this quadratic is convex, the point is also the global minimizer.

\end{workedexamplebox}

%%%%%%%%%%%%%%%%%%%%%%%%%%%%%%%%%%%%%%%%%%%%%%%%%%%%%%%%%%%%
\subsection{Quadratic Functions}
\label{subsec:quadratic-objectives}
%%%%%%%%%%%%%%%%%%%%%%%%%%%%%%%%%%%%%%%%%%%%%%%%%%%%%%%%%%%%

A general quadratic objective can be written

\[
\boxed{
f(\mathbf{x})
=
\frac12
\mathbf{x}^TQ\mathbf{x}
+
\mathbf{c}^T\mathbf{x}
+
d.
}
\]

If \(Q\) is symmetric, then

\[
\boxed{
\nabla f
=
Q\mathbf{x}
+
\mathbf{c},
}
\]

and

\[
\boxed{
\nabla^2f
=
Q.
}
\]

%%%%%%%%%%%%%%%%%%%%%%%%%%%%%%%%%%%%%%%%%%%%%%%%%%%%%%%%%%%%
\subsection{Quadratic Minimizer}
\label{subsec:quadratic-minimizer}
%%%%%%%%%%%%%%%%%%%%%%%%%%%%%%%%%%%%%%%%%%%%%%%%%%%%%%%%%%%%

If

\[
Q\succ0,
\]

the unique unconstrained minimizer satisfies

\[
Q\mathbf{x}+\mathbf{c}=0.
\]

Thus,

\[
\boxed{
\mathbf{x}^*
=
-Q^{-1}\mathbf{c}.
}
\]

Positive definiteness guarantees uniqueness.

%%%%%%%%%%%%%%%%%%%%%%%%%%%%%%%%%%%%%%%%%%%%%%%%%%%%%%%%%%%%
\subsection{Least-Squares Optimization}
\label{subsec:least-squares-optimization}
%%%%%%%%%%%%%%%%%%%%%%%%%%%%%%%%%%%%%%%%%%%%%%%%%%%%%%%%%%%%

Given

\[
A\in\mathbb{R}^{m\times n}
\]

and

\[
\mathbf{b}\in\mathbb{R}^m,
\]

least squares solves

\[
\boxed{
\min_{\mathbf{x}}
\frac12
\|
A\mathbf{x}-\mathbf{b}
\|_2^2.
}
\]

The gradient is

\[
\boxed{
\nabla f
=
A^T
(
A\mathbf{x}-\mathbf{b}
).
}
\]

%%%%%%%%%%%%%%%%%%%%%%%%%%%%%%%%%%%%%%%%%%%%%%%%%%%%%%%%%%%%
\subsection{Normal Equations}
\label{subsec:normal-equations-optimization}
%%%%%%%%%%%%%%%%%%%%%%%%%%%%%%%%%%%%%%%%%%%%%%%%%%%%%%%%%%%%

The first-order condition gives

\[
A^T
(
A\mathbf{x}-\mathbf{b}
)
=
0.
\]

Thus,

\[
\boxed{
A^TA\mathbf{x}
=
A^T\mathbf{b}.
}
\]

These are the normal equations.

In numerical computation, QR or SVD methods are often preferable to
forming \(A^TA\) explicitly.

%%%%%%%%%%%%%%%%%%%%%%%%%%%%%%%%%%%%%%%%%%%%%%%%%%%%%%%%%%%%
\subsection{Existence of an Optimum}
\label{subsec:existence-optimum}
%%%%%%%%%%%%%%%%%%%%%%%%%%%%%%%%%%%%%%%%%%%%%%%%%%%%%%%%%%%%

Optimization conditions are useful only if a minimizer exists.

One standard result is the extreme-value theorem.

If \(f\) is continuous and the feasible set \(\mathcal{F}\) is nonempty
and compact, then \(f\) attains both a minimum and a maximum on
\(\mathcal{F}\).

%%%%%%%%%%%%%%%%%%%%%%%%%%%%%%%%%%%%%%%%%%%%%%%%%%%%%%%%%%%%
\subsection{Coercive Functions}
\label{subsec:coercive-functions}
%%%%%%%%%%%%%%%%%%%%%%%%%%%%%%%%%%%%%%%%%%%%%%%%%%%%%%%%%%%%

A function is coercive if

\[
\boxed{
f(\mathbf{x})
\to\infty
\qquad
\text{as }
\|\mathbf{x}\|\to\infty.
}
\]

A continuous coercive function on \(\mathbb{R}^n\) attains a global
minimum.

Coercivity prevents minimizing sequences from escaping to infinity.

%%%%%%%%%%%%%%%%%%%%%%%%%%%%%%%%%%%%%%%%%%%%%%%%%%%%%%%%%%%%
\subsection{Convex Sets}
\label{subsec:convex-sets-optimization}
%%%%%%%%%%%%%%%%%%%%%%%%%%%%%%%%%%%%%%%%%%%%%%%%%%%%%%%%%%%%

\begin{definitionbox}{Convex Set}
A set \(C\subseteq\mathbb{R}^n\) is convex if, for every
\(\mathbf{x},\mathbf{y}\in C\) and every \(0\leq\theta\leq1\),

\[
\boxed{
\theta\mathbf{x}
+
(1-\theta)\mathbf{y}
\in C.
}
\]
\end{definitionbox}

Geometrically, every line segment connecting two points of \(C\) remains
inside \(C\).

%%%%%%%%%%%%%%%%%%%%%%%%%%%%%%%%%%%%%%%%%%%%%%%%%%%%%%%%%%%%
\subsection{Examples of Convex Sets}
\label{subsec:examples-convex-sets}
%%%%%%%%%%%%%%%%%%%%%%%%%%%%%%%%%%%%%%%%%%%%%%%%%%%%%%%%%%%%

Examples include:

\[
\boxed{
\text{intervals},
}
\]

\[
\boxed{
\text{affine subspaces},
}
\]

\[
\boxed{
\text{half-spaces},
}
\]

\[
\boxed{
\text{norm balls},
}
\]

and

\[
\boxed{
\text{intersections of convex sets}.
}
\]

%%%%%%%%%%%%%%%%%%%%%%%%%%%%%%%%%%%%%%%%%%%%%%%%%%%%%%%%%%%%
\subsection{Convex Functions}
\label{subsec:convex-functions-optimization}
%%%%%%%%%%%%%%%%%%%%%%%%%%%%%%%%%%%%%%%%%%%%%%%%%%%%%%%%%%%%

\begin{definitionbox}{Convex Function}
A function \(f:C\to\mathbb{R}\) on a convex set \(C\) is convex if

\[
\boxed{
f
(
\theta\mathbf{x}
+
(1-\theta)\mathbf{y}
)
\leq
\theta f(\mathbf{x})
+
(1-\theta)f(\mathbf{y})
}
\]

for all \(\mathbf{x},\mathbf{y}\in C\) and \(0\leq\theta\leq1\).
\end{definitionbox}

%%%%%%%%%%%%%%%%%%%%%%%%%%%%%%%%%%%%%%%%%%%%%%%%%%%%%%%%%%%%
\subsection{Strict Convexity}
\label{subsec:strict-convexity}
%%%%%%%%%%%%%%%%%%%%%%%%%%%%%%%%%%%%%%%%%%%%%%%%%%%%%%%%%%%%

A function is strictly convex if

\[
\boxed{
f
(
\theta\mathbf{x}
+
(1-\theta)\mathbf{y}
)
<
\theta f(\mathbf{x})
+
(1-\theta)f(\mathbf{y})
}
\]

for distinct \(\mathbf{x},\mathbf{y}\) and

\[
0<\theta<1.
\]

A strictly convex function has at most one global minimizer on a convex
feasible set.

%%%%%%%%%%%%%%%%%%%%%%%%%%%%%%%%%%%%%%%%%%%%%%%%%%%%%%%%%%%%
\subsection{First-Order Characterization of Convexity}
\label{subsec:first-order-convexity}
%%%%%%%%%%%%%%%%%%%%%%%%%%%%%%%%%%%%%%%%%%%%%%%%%%%%%%%%%%%%

If \(f\) is differentiable and convex, then

\[
\boxed{
f(\mathbf{y})
\geq
f(\mathbf{x})
+
\nabla f(\mathbf{x})^T
(
\mathbf{y}-\mathbf{x}
)
}
\]

for all \(\mathbf{x},\mathbf{y}\) in the convex domain.

Thus the tangent hyperplane lies below the graph.

%%%%%%%%%%%%%%%%%%%%%%%%%%%%%%%%%%%%%%%%%%%%%%%%%%%%%%%%%%%%
\subsection{Second-Order Convexity Test}
\label{subsec:second-order-convexity}
%%%%%%%%%%%%%%%%%%%%%%%%%%%%%%%%%%%%%%%%%%%%%%%%%%%%%%%%%%%%

If \(f\) is twice differentiable on a convex domain, then

\[
\boxed{
\nabla^2f(\mathbf{x})
\succeq0
}
\]

for every \(\mathbf{x}\) implies convexity.

If

\[
\boxed{
\nabla^2f(\mathbf{x})
\succ0
}
\]

everywhere, then \(f\) is strictly convex.

%%%%%%%%%%%%%%%%%%%%%%%%%%%%%%%%%%%%%%%%%%%%%%%%%%%%%%%%%%%%
\subsection{Why Convexity Matters}
\label{subsec:why-convexity-matters}
%%%%%%%%%%%%%%%%%%%%%%%%%%%%%%%%%%%%%%%%%%%%%%%%%%%%%%%%%%%%

For a convex optimization problem:

\[
\boxed{
\text{every local minimizer is global}.
}
\]

If the objective is strictly convex, a minimizer is unique when it
exists.

Therefore convexity transforms a difficult global problem into one where
local information can often determine the global solution.

%%%%%%%%%%%%%%%%%%%%%%%%%%%%%%%%%%%%%%%%%%%%%%%%%%%%%%%%%%%%
\subsection{Convex Optimization Problem}
\label{subsec:convex-optimization-problem}
%%%%%%%%%%%%%%%%%%%%%%%%%%%%%%%%%%%%%%%%%%%%%%%%%%%%%%%%%%%%

A standard convex problem has:

\[
\boxed{
\text{convex objective }f,
}
\]

\[
\boxed{
\text{convex inequality functions }g_i,
}
\]

and

\[
\boxed{
\text{affine equality constraints}.
}
\]

It can be written

\[
\boxed{
\begin{aligned}
\min_{\mathbf{x}}
\quad&
f(\mathbf{x})
\\
\text{subject to}
\quad&
g_i(\mathbf{x})\leq0,
\\
&
A\mathbf{x}=\mathbf{b}.
\end{aligned}
}
\]

%%%%%%%%%%%%%%%%%%%%%%%%%%%%%%%%%%%%%%%%%%%%%%%%%%%%%%%%%%%%
\subsection{Equality-Constrained Optimization}
\label{subsec:equality-constrained-optimization}
%%%%%%%%%%%%%%%%%%%%%%%%%%%%%%%%%%%%%%%%%%%%%%%%%%%%%%%%%%%%

Consider

\[
\boxed{
\min_{\mathbf{x}}
f(\mathbf{x})
}
\]

subject to

\[
\boxed{
h_j(\mathbf{x})=0,
\qquad
j=1,\ldots,m.
}
\]

At a regular constrained optimum, feasible movement is restricted to
directions tangent to the constraint surface.

%%%%%%%%%%%%%%%%%%%%%%%%%%%%%%%%%%%%%%%%%%%%%%%%%%%%%%%%%%%%
\subsection{Lagrangian}
\label{subsec:optimization-lagrangian}
%%%%%%%%%%%%%%%%%%%%%%%%%%%%%%%%%%%%%%%%%%%%%%%%%%%%%%%%%%%%

For equality constraints, define the Lagrangian

\[
\boxed{
\mathcal{L}
(
\mathbf{x},
\boldsymbol{\lambda}
)
=
f(\mathbf{x})
+
\sum_{j=1}^{m}
\lambda_jh_j(\mathbf{x}).
}
\]

The variables \(\lambda_j\) are Lagrange multipliers.

%%%%%%%%%%%%%%%%%%%%%%%%%%%%%%%%%%%%%%%%%%%%%%%%%%%%%%%%%%%%
\subsection{Lagrange Multiplier Conditions}
\label{subsec:lagrange-multiplier-conditions}
%%%%%%%%%%%%%%%%%%%%%%%%%%%%%%%%%%%%%%%%%%%%%%%%%%%%%%%%%%%%

Under suitable regularity conditions, a local constrained optimum
satisfies

\[
\boxed{
\nabla_{\mathbf{x}}
\mathcal{L}
=
\mathbf{0},
}
\]

together with

\[
\boxed{
h_j(\mathbf{x})=0.
}
\]

Equivalently,

\[
\boxed{
\nabla f(\mathbf{x}^*)
+
\sum_{j=1}^{m}
\lambda_j
\nabla h_j(\mathbf{x}^*)
=
0.
}
\]

%%%%%%%%%%%%%%%%%%%%%%%%%%%%%%%%%%%%%%%%%%%%%%%%%%%%%%%%%%%%
\subsection{Geometric Interpretation of Lagrange Multipliers}
\label{subsec:lagrange-geometric-interpretation}
%%%%%%%%%%%%%%%%%%%%%%%%%%%%%%%%%%%%%%%%%%%%%%%%%%%%%%%%%%%%

For one constraint

\[
h(\mathbf{x})=0,
\]

the condition

\[
\nabla f
+
\lambda\nabla h
=
0
\]

means

\[
\boxed{
\nabla f
\parallel
\nabla h.
}
\]

The objective gradient is normal to the feasible constraint surface at
the optimum.

%%%%%%%%%%%%%%%%%%%%%%%%%%%%%%%%%%%%%%%%%%%%%%%%%%%%%%%%%%%%
\subsection{Worked Example: Equality Constraint}
\label{subsec:worked-lagrange-multiplier}
%%%%%%%%%%%%%%%%%%%%%%%%%%%%%%%%%%%%%%%%%%%%%%%%%%%%%%%%%%%%

\begin{workedexamplebox}{Minimizing Distance on a Line}

Minimize

\[
f(x,y)
=
x^2+y^2
\]

subject to

\[
x+y=1.
\]

Define

\[
\mathcal{L}
=
x^2+y^2
+
\lambda(x+y-1).
\]

The first-order equations are

\[
2x+\lambda=0,
\]

\[
2y+\lambda=0,
\]

and

\[
x+y=1.
\]

The first two equations imply

\[
x=y.
\]

Therefore,

\[
2x=1,
\]

so

\[
x=y=\frac12.
\]

Hence,

\begin{answerbox}
\[
\boxed{
(x^*,y^*)
=
\left(
\frac12,
\frac12
\right).
}
\]
\end{answerbox}

\end{workedexamplebox}

%%%%%%%%%%%%%%%%%%%%%%%%%%%%%%%%%%%%%%%%%%%%%%%%%%%%%%%%%%%%
\subsection{Inequality Constraints}
\label{subsec:inequality-constraints-optimization}
%%%%%%%%%%%%%%%%%%%%%%%%%%%%%%%%%%%%%%%%%%%%%%%%%%%%%%%%%%%%

Consider

\[
\boxed{
g_i(\mathbf{x})\leq0.
}
\]

At a feasible point, an inequality is active if

\[
\boxed{
g_i(\mathbf{x})=0.
}
\]

It is inactive if

\[
\boxed{
g_i(\mathbf{x})<0.
}
\]

Only active inequalities locally restrict feasible movement.

%%%%%%%%%%%%%%%%%%%%%%%%%%%%%%%%%%%%%%%%%%%%%%%%%%%%%%%%%%%%
\subsection{Karush--Kuhn--Tucker Conditions}
\label{subsec:kkt-conditions}
%%%%%%%%%%%%%%%%%%%%%%%%%%%%%%%%%%%%%%%%%%%%%%%%%%%%%%%%%%%%

Consider

\[
\begin{aligned}
\min_{\mathbf{x}}
\quad&
f(\mathbf{x})
\\
\text{subject to}
\quad&
g_i(\mathbf{x})\leq0,
\\
&
h_j(\mathbf{x})=0.
\end{aligned}
\]

Define

\[
\boxed{
\mathcal{L}
=
f(\mathbf{x})
+
\sum_i
\mu_i g_i(\mathbf{x})
+
\sum_j
\lambda_j h_j(\mathbf{x}).
}
\]

Under an appropriate constraint qualification, a local optimum satisfies
the KKT conditions.

%%%%%%%%%%%%%%%%%%%%%%%%%%%%%%%%%%%%%%%%%%%%%%%%%%%%%%%%%%%%
\subsection{KKT Stationarity}
\label{subsec:kkt-stationarity}
%%%%%%%%%%%%%%%%%%%%%%%%%%%%%%%%%%%%%%%%%%%%%%%%%%%%%%%%%%%%

Stationarity requires

\[
\boxed{
\nabla f(\mathbf{x}^*)
+
\sum_i
\mu_i
\nabla g_i(\mathbf{x}^*)
+
\sum_j
\lambda_j
\nabla h_j(\mathbf{x}^*)
=
0.
}
\]

%%%%%%%%%%%%%%%%%%%%%%%%%%%%%%%%%%%%%%%%%%%%%%%%%%%%%%%%%%%%
\subsection{Primal Feasibility}
\label{subsec:kkt-primal-feasibility}
%%%%%%%%%%%%%%%%%%%%%%%%%%%%%%%%%%%%%%%%%%%%%%%%%%%%%%%%%%%%

The candidate solution must satisfy the original constraints:

\[
\boxed{
g_i(\mathbf{x}^*)\leq0,
}
\]

and

\[
\boxed{
h_j(\mathbf{x}^*)=0.
}
\]

%%%%%%%%%%%%%%%%%%%%%%%%%%%%%%%%%%%%%%%%%%%%%%%%%%%%%%%%%%%%
\subsection{Dual Feasibility}
\label{subsec:kkt-dual-feasibility}
%%%%%%%%%%%%%%%%%%%%%%%%%%%%%%%%%%%%%%%%%%%%%%%%%%%%%%%%%%%%

For the convention

\[
g_i(\mathbf{x})\leq0,
\]

the inequality multipliers satisfy

\[
\boxed{
\mu_i\geq0.
}
\]

The equality multipliers \(\lambda_j\) are unrestricted in sign.

%%%%%%%%%%%%%%%%%%%%%%%%%%%%%%%%%%%%%%%%%%%%%%%%%%%%%%%%%%%%
\subsection{Complementary Slackness}
\label{subsec:complementary-slackness}
%%%%%%%%%%%%%%%%%%%%%%%%%%%%%%%%%%%%%%%%%%%%%%%%%%%%%%%%%%%%

For every inequality constraint,

\[
\boxed{
\mu_i
g_i(\mathbf{x}^*)
=
0.
}
\]

Therefore:

\[
g_i(\mathbf{x}^*)<0
\Rightarrow
\mu_i=0.
\]

An inactive constraint has zero multiplier.

%%%%%%%%%%%%%%%%%%%%%%%%%%%%%%%%%%%%%%%%%%%%%%%%%%%%%%%%%%%%
\subsection{KKT Conditions Summary}
\label{subsec:kkt-summary}
%%%%%%%%%%%%%%%%%%%%%%%%%%%%%%%%%%%%%%%%%%%%%%%%%%%%%%%%%%%%

The KKT conditions are:

\[
\boxed{
\text{stationarity},
}
\]

\[
\boxed{
\text{primal feasibility},
}
\]

\[
\boxed{
\text{dual feasibility},
}
\]

and

\[
\boxed{
\text{complementary slackness}.
}
\]

For convex problems satisfying appropriate regularity conditions, these
conditions are often sufficient for global optimality.

%%%%%%%%%%%%%%%%%%%%%%%%%%%%%%%%%%%%%%%%%%%%%%%%%%%%%%%%%%%%
\subsection{Worked Example: Simple KKT Problem}
\label{subsec:worked-kkt-problem}
%%%%%%%%%%%%%%%%%%%%%%%%%%%%%%%%%%%%%%%%%%%%%%%%%%%%%%%%%%%%

\begin{workedexamplebox}{Bound-Constrained Minimum}

Minimize

\[
f(x)
=
(x-3)^2
\]

subject to

\[
x\leq1.
\]

Write

\[
g(x)=x-1\leq0.
\]

The Lagrangian is

\[
\mathcal{L}
=
(x-3)^2
+
\mu(x-1).
\]

Stationarity gives

\[
2(x-3)+\mu=0.
\]

The unconstrained minimizer is \(x=3\), which is infeasible.

Therefore the constraint must be active:

\[
x=1.
\]

Then

\[
2(1-3)+\mu=0,
\]

so

\[
\mu=4.
\]

Since

\[
\mu\geq0,
\]

the KKT conditions hold.

\begin{answerbox}
\[
\boxed{
x^*=1,
\qquad
\mu^*=4.
}
\]
\end{answerbox}

\end{workedexamplebox}

%%%%%%%%%%%%%%%%%%%%%%%%%%%%%%%%%%%%%%%%%%%%%%%%%%%%%%%%%%%%
\subsection{Constraint Qualifications}
\label{subsec:constraint-qualifications}
%%%%%%%%%%%%%%%%%%%%%%%%%%%%%%%%%%%%%%%%%%%%%%%%%%%%%%%%%%%%

KKT conditions are not automatically necessary for every constrained
problem.

A constraint qualification is a regularity condition ensuring that the
constraints behave well enough for multiplier conditions to apply.

Examples include:

\[
\boxed{
\text{LICQ},
}
\]

\[
\boxed{
\text{MFCQ},
}
\]

and, for convex problems,

\[
\boxed{
\text{Slater's condition}.
}
\]

These are developed more deeply in nonlinear and convex optimization.

%%%%%%%%%%%%%%%%%%%%%%%%%%%%%%%%%%%%%%%%%%%%%%%%%%%%%%%%%%%%
\subsection{Lagrange Multipliers as Sensitivities}
\label{subsec:lagrange-multiplier-sensitivity}
%%%%%%%%%%%%%%%%%%%%%%%%%%%%%%%%%%%%%%%%%%%%%%%%%%%%%%%%%%%%

Suppose a constraint depends on a resource parameter \(b\).

Under appropriate differentiability and regularity,

\[
\boxed{
\lambda^*
}
\]

can measure how the optimal objective value changes when \(b\) changes.

This motivates the economic interpretation of multipliers as

\[
\boxed{
\text{shadow prices}.
}
\]

The sign depends on the exact constraint and Lagrangian convention.

%%%%%%%%%%%%%%%%%%%%%%%%%%%%%%%%%%%%%%%%%%%%%%%%%%%%%%%%%%%%
\subsection{Regularization}
\label{subsec:optimization-regularization}
%%%%%%%%%%%%%%%%%%%%%%%%%%%%%%%%%%%%%%%%%%%%%%%%%%%%%%%%%%%%

Regularization modifies an objective by adding a penalty term:

\[
\boxed{
\min_{\mathbf{x}}
f(\mathbf{x})
+
\lambda R(\mathbf{x}).
}
\]

Here:

\[
R(\mathbf{x})
=
\text{regularizer},
\]

and

\[
\lambda\geq0
\]

controls penalty strength.

%%%%%%%%%%%%%%%%%%%%%%%%%%%%%%%%%%%%%%%%%%%%%%%%%%%%%%%%%%%%
\subsection{\(\ell_2\) Regularization}
\label{subsec:l2-regularization}
%%%%%%%%%%%%%%%%%%%%%%%%%%%%%%%%%%%%%%%%%%%%%%%%%%%%%%%%%%%%

A common regularizer is

\[
\boxed{
R(\mathbf{x})
=
\frac12
\|\mathbf{x}\|_2^2.
}
\]

Then

\[
\boxed{
\min_{\mathbf{x}}
f(\mathbf{x})
+
\frac{\lambda}{2}
\|\mathbf{x}\|_2^2.
}
\]

This discourages large coefficient magnitude and can improve conditioning.

%%%%%%%%%%%%%%%%%%%%%%%%%%%%%%%%%%%%%%%%%%%%%%%%%%%%%%%%%%%%
\subsection{\(\ell_1\) Regularization}
\label{subsec:l1-regularization}
%%%%%%%%%%%%%%%%%%%%%%%%%%%%%%%%%%%%%%%%%%%%%%%%%%%%%%%%%%%%

Another common regularizer is

\[
\boxed{
R(\mathbf{x})
=
\|\mathbf{x}\|_1
=
\sum_i|x_i|.
}
\]

The resulting problem

\[
\boxed{
\min_{\mathbf{x}}
f(\mathbf{x})
+
\lambda
\|\mathbf{x}\|_1
}
\]

can encourage sparse solutions.

The objective is generally nondifferentiable when a component equals
zero.

%%%%%%%%%%%%%%%%%%%%%%%%%%%%%%%%%%%%%%%%%%%%%%%%%%%%%%%%%%%%
\subsection{Penalty Methods}
\label{subsec:penalty-methods-preview}
%%%%%%%%%%%%%%%%%%%%%%%%%%%%%%%%%%%%%%%%%%%%%%%%%%%%%%%%%%%%

Constraints can sometimes be incorporated into the objective using
penalties.

For an equality constraint

\[
h(\mathbf{x})=0,
\]

one may consider

\[
\boxed{
f_\rho(\mathbf{x})
=
f(\mathbf{x})
+
\frac{\rho}{2}
\|
h(\mathbf{x})
\|^2.
}
\]

Large \(\rho\) penalizes constraint violation.

Penalty methods do not generally enforce exact feasibility for finite
\(\rho\).

%%%%%%%%%%%%%%%%%%%%%%%%%%%%%%%%%%%%%%%%%%%%%%%%%%%%%%%%%%%%
\subsection{Barrier Methods Preview}
\label{subsec:barrier-methods-preview}
%%%%%%%%%%%%%%%%%%%%%%%%%%%%%%%%%%%%%%%%%%%%%%%%%%%%%%%%%%%%

For an inequality

\[
g_i(\mathbf{x})<0,
\]

a logarithmic barrier uses

\[
\boxed{
-\mu
\log
[
-g_i(\mathbf{x})
].
}
\]

Thus a constrained problem can be approximated by

\[
\boxed{
\min_{\mathbf{x}}
f(\mathbf{x})
-
\mu
\sum_i
\log
[
-g_i(\mathbf{x})
].
}
\]

Barrier terms become large near the boundary.

%%%%%%%%%%%%%%%%%%%%%%%%%%%%%%%%%%%%%%%%%%%%%%%%%%%%%%%%%%%%
\subsection{Gradient Descent}
\label{subsec:gradient-descent}
%%%%%%%%%%%%%%%%%%%%%%%%%%%%%%%%%%%%%%%%%%%%%%%%%%%%%%%%%%%%

Gradient descent updates

\[
\boxed{
\mathbf{x}_{k+1}
=
\mathbf{x}_k
-
\alpha_k
\nabla f(\mathbf{x}_k),
}
\]

where

\[
\alpha_k>0
\]

is the step size.

Because

\[
-\nabla f
\]

is a local descent direction, sufficiently small steps can reduce the
objective when the gradient is nonzero.

%%%%%%%%%%%%%%%%%%%%%%%%%%%%%%%%%%%%%%%%%%%%%%%%%%%%%%%%%%%%
\subsection{Worked Example: Gradient Descent Step}
\label{subsec:worked-gradient-descent}
%%%%%%%%%%%%%%%%%%%%%%%%%%%%%%%%%%%%%%%%%%%%%%%%%%%%%%%%%%%%

\begin{workedexamplebox}{One Gradient Step}

Let

\[
f(x)
=
(x-4)^2.
\]

Then

\[
f'(x)
=
2(x-4).
\]

Start at

\[
x_0=0
\]

with

\[
\alpha=0.1.
\]

The gradient is

\[
f'(0)=-8.
\]

Therefore,

\[
x_1
=
0
-
0.1(-8).
\]

Hence,

\begin{answerbox}
\[
\boxed{
x_1=0.8.
}
\]
\end{answerbox}

\end{workedexamplebox}

%%%%%%%%%%%%%%%%%%%%%%%%%%%%%%%%%%%%%%%%%%%%%%%%%%%%%%%%%%%%
\subsection{Step Size}
\label{subsec:optimization-step-size}
%%%%%%%%%%%%%%%%%%%%%%%%%%%%%%%%%%%%%%%%%%%%%%%%%%%%%%%%%%%%

The step size controls movement along a search direction.

If it is too small:

\[
\boxed{
\text{convergence can be slow}.
}
\]

If it is too large:

\[
\boxed{
\text{the method may overshoot or diverge}.
}
\]

Choosing step sizes is therefore a major part of numerical optimization.

%%%%%%%%%%%%%%%%%%%%%%%%%%%%%%%%%%%%%%%%%%%%%%%%%%%%%%%%%%%%
\subsection{Line Search}
\label{subsec:line-search-preview}
%%%%%%%%%%%%%%%%%%%%%%%%%%%%%%%%%%%%%%%%%%%%%%%%%%%%%%%%%%%%

Given a search direction \(\mathbf{p}_k\), a line-search method chooses

\[
\alpha_k
\]

and updates

\[
\boxed{
\mathbf{x}_{k+1}
=
\mathbf{x}_k
+
\alpha_k\mathbf{p}_k.
}
\]

The one-dimensional function

\[
\boxed{
\phi(\alpha)
=
f
(
\mathbf{x}_k+\alpha\mathbf{p}_k
)
}
\]

is examined to find a useful step.

%%%%%%%%%%%%%%%%%%%%%%%%%%%%%%%%%%%%%%%%%%%%%%%%%%%%%%%%%%%%
\subsection{Newton's Method for Optimization}
\label{subsec:newtons-method-optimization}
%%%%%%%%%%%%%%%%%%%%%%%%%%%%%%%%%%%%%%%%%%%%%%%%%%%%%%%%%%%%

Near a point \(\mathbf{x}_k\), approximate

\[
f(\mathbf{x}_k+\mathbf{p})
\]

by its quadratic Taylor expansion:

\[
f(\mathbf{x}_k)
+
\nabla f(\mathbf{x}_k)^T\mathbf{p}
+
\frac12
\mathbf{p}^T
\nabla^2f(\mathbf{x}_k)
\mathbf{p}.
\]

Minimizing this model gives the Newton equation

\[
\boxed{
\nabla^2f(\mathbf{x}_k)\mathbf{p}_k
=
-
\nabla f(\mathbf{x}_k).
}
\]

Then

\[
\boxed{
\mathbf{x}_{k+1}
=
\mathbf{x}_k+\mathbf{p}_k.
}
\]

%%%%%%%%%%%%%%%%%%%%%%%%%%%%%%%%%%%%%%%%%%%%%%%%%%%%%%%%%%%%
\subsection{Newton Step}
\label{subsec:newton-step-optimization}
%%%%%%%%%%%%%%%%%%%%%%%%%%%%%%%%%%%%%%%%%%%%%%%%%%%%%%%%%%%%

If the Hessian is invertible,

\[
\boxed{
\mathbf{p}_k
=
-
[
\nabla^2f(\mathbf{x}_k)
]^{-1}
\nabla f(\mathbf{x}_k).
}
\]

In numerical computation, one solves the linear system rather than
forming the matrix inverse explicitly.

%%%%%%%%%%%%%%%%%%%%%%%%%%%%%%%%%%%%%%%%%%%%%%%%%%%%%%%%%%%%
\subsection{Worked Example: Newton Method}
\label{subsec:worked-newton-optimization}
%%%%%%%%%%%%%%%%%%%%%%%%%%%%%%%%%%%%%%%%%%%%%%%%%%%%%%%%%%%%

\begin{workedexamplebox}{Newton Step for a Quadratic}

Let

\[
f(x)
=
x^2-6x+10.
\]

Then

\[
f'(x)=2x-6,
\]

and

\[
f''(x)=2.
\]

Newton's update is

\[
x_{k+1}
=
x_k
-
\frac{
2x_k-6
}{
2
}.
\]

Thus,

\[
x_{k+1}
=
3.
\]

For this quadratic, Newton's method reaches the minimizer in one step from
any starting point.

\begin{answerbox}
\[
\boxed{
x^*=3.
}
\]
\end{answerbox}

\end{workedexamplebox}

%%%%%%%%%%%%%%%%%%%%%%%%%%%%%%%%%%%%%%%%%%%%%%%%%%%%%%%%%%%%
\subsection{Trust-Region Methods Preview}
\label{subsec:trust-region-preview}
%%%%%%%%%%%%%%%%%%%%%%%%%%%%%%%%%%%%%%%%%%%%%%%%%%%%%%%%%%%%

Instead of choosing a step length along one direction, a trust-region
method minimizes a local model subject to

\[
\boxed{
\|\mathbf{p}\|
\leq
\Delta_k.
}
\]

The radius \(\Delta_k\) represents the region where the local model is
trusted.

The radius is adjusted based on agreement between predicted and actual
objective reduction.

%%%%%%%%%%%%%%%%%%%%%%%%%%%%%%%%%%%%%%%%%%%%%%%%%%%%%%%%%%%%
\subsection{First-Order and Second-Order Methods}
\label{subsec:first-second-order-methods}
%%%%%%%%%%%%%%%%%%%%%%%%%%%%%%%%%%%%%%%%%%%%%%%%%%%%%%%%%%%%

First-order methods use gradients:

\[
\boxed{
\nabla f.
}
\]

Second-order methods use Hessian information:

\[
\boxed{
\nabla^2f.
}
\]

First-order methods are often cheaper per iteration.

Second-order methods can converge much faster near a solution when the
Hessian is informative and computationally manageable.

%%%%%%%%%%%%%%%%%%%%%%%%%%%%%%%%%%%%%%%%%%%%%%%%%%%%%%%%%%%%
\subsection{Quasi-Newton Methods Preview}
\label{subsec:quasi-newton-preview}
%%%%%%%%%%%%%%%%%%%%%%%%%%%%%%%%%%%%%%%%%%%%%%%%%%%%%%%%%%%%

Quasi-Newton methods approximate the Hessian or its inverse using
successive gradient information.

They avoid computing exact second derivatives.

Important examples include:

\[
\boxed{
\text{BFGS},
}
\]

and

\[
\boxed{
\text{L-BFGS}.
}
\]

These methods are developed more deeply in numerical optimization.

%%%%%%%%%%%%%%%%%%%%%%%%%%%%%%%%%%%%%%%%%%%%%%%%%%%%%%%%%%%%
\subsection{Optimality Gap}
\label{subsec:optimality-gap}
%%%%%%%%%%%%%%%%%%%%%%%%%%%%%%%%%%%%%%%%%%%%%%%%%%%%%%%%%%%%

If a known lower bound \(L\) satisfies

\[
L\leq f^*,
\]

and a feasible point has value

\[
f(\mathbf{x}),
\]

then

\[
\boxed{
f(\mathbf{x})-L
}
\]

is an upper bound on the distance from the current objective value to the
global optimum.

Optimization algorithms often use such gaps as stopping criteria.

%%%%%%%%%%%%%%%%%%%%%%%%%%%%%%%%%%%%%%%%%%%%%%%%%%%%%%%%%%%%
\subsection{Stopping Criteria}
\label{subsec:optimization-stopping-criteria}
%%%%%%%%%%%%%%%%%%%%%%%%%%%%%%%%%%%%%%%%%%%%%%%%%%%%%%%%%%%%

An iterative method may stop when one or more conditions hold:

\[
\boxed{
\|\nabla f(\mathbf{x}_k)\|
\leq\varepsilon,
}
\]

\[
\boxed{
\|
\mathbf{x}_{k+1}-\mathbf{x}_k
\|
\leq\varepsilon,
}
\]

or

\[
\boxed{
|
f(\mathbf{x}_{k+1})
-
f(\mathbf{x}_k)
|
\leq\varepsilon.
}
\]

For constrained problems, feasibility residuals must also be checked.

%%%%%%%%%%%%%%%%%%%%%%%%%%%%%%%%%%%%%%%%%%%%%%%%%%%%%%%%%%%%
\subsection{Scaling}
\label{subsec:optimization-scaling}
%%%%%%%%%%%%%%%%%%%%%%%%%%%%%%%%%%%%%%%%%%%%%%%%%%%%%%%%%%%%

Suppose

\[
x_1
\]

naturally has magnitude

\[
10^{-6},
\]

while

\[
x_2
\]

has magnitude

\[
10^6.
\]

Poor scaling can cause numerical difficulties.

Rescaling variables and constraints can improve:

\[
\boxed{
\text{conditioning},
}
\]

\[
\boxed{
\text{step selection},
}
\]

and

\[
\boxed{
\text{solver reliability}.
}
\]

%%%%%%%%%%%%%%%%%%%%%%%%%%%%%%%%%%%%%%%%%%%%%%%%%%%%%%%%%%%%
\subsection{Conditioning}
\label{subsec:optimization-conditioning}
%%%%%%%%%%%%%%%%%%%%%%%%%%%%%%%%%%%%%%%%%%%%%%%%%%%%%%%%%%%%

A problem is ill conditioned if small changes in the data can produce
large changes in the solution.

For a quadratic objective,

\[
f(x)
=
\frac12x^TQx-c^Tx,
\]

the eigenvalue spread of \(Q\) influences numerical difficulty.

A large condition number can cause slow convergence for gradient-based
methods.

%%%%%%%%%%%%%%%%%%%%%%%%%%%%%%%%%%%%%%%%%%%%%%%%%%%%%%%%%%%%
\subsection{Deterministic Optimization}
\label{subsec:deterministic-optimization}
%%%%%%%%%%%%%%%%%%%%%%%%%%%%%%%%%%%%%%%%%%%%%%%%%%%%%%%%%%%%

In deterministic optimization, the model data and objective are treated
as known.

A problem may have form

\[
\boxed{
\min_x f(x).
}
\]

Repeated evaluation at the same \(x\) gives the same mathematical
objective value.

%%%%%%%%%%%%%%%%%%%%%%%%%%%%%%%%%%%%%%%%%%%%%%%%%%%%%%%%%%%%
\subsection{Stochastic Optimization Preview}
\label{subsec:stochastic-optimization-preview}
%%%%%%%%%%%%%%%%%%%%%%%%%%%%%%%%%%%%%%%%%%%%%%%%%%%%%%%%%%%%

In stochastic optimization, uncertainty enters the model.

A common form is

\[
\boxed{
\min_x
\mathbb{E}
[
F(x,\xi)
],
}
\]

where

\[
\xi
\]

is random.

Another form includes probabilistic constraints.

Stochastic optimization is developed later in Chapter 10.

%%%%%%%%%%%%%%%%%%%%%%%%%%%%%%%%%%%%%%%%%%%%%%%%%%%%%%%%%%%%
\subsection{Multiobjective Optimization Preview}
\label{subsec:multiobjective-optimization-preview}
%%%%%%%%%%%%%%%%%%%%%%%%%%%%%%%%%%%%%%%%%%%%%%%%%%%%%%%%%%%%

Some problems have several competing objectives:

\[
\boxed{
\min_x
\begin{pmatrix}
f_1(x)\\
f_2(x)\\
\vdots\\
f_k(x)
\end{pmatrix}.
}
\]

There may be no single point minimizing every objective simultaneously.

Instead, one studies tradeoffs and Pareto optimality.

%%%%%%%%%%%%%%%%%%%%%%%%%%%%%%%%%%%%%%%%%%%%%%%%%%%%%%%%%%%%
\subsection{Discrete Versus Continuous Optimization}
\label{subsec:discrete-continuous-optimization}
%%%%%%%%%%%%%%%%%%%%%%%%%%%%%%%%%%%%%%%%%%%%%%%%%%%%%%%%%%%%

Continuous optimization uses variables such as

\[
\boxed{
x\in\mathbb{R}^n.
}
\]

Discrete optimization restricts some or all variables to sets such as

\[
\boxed{
\mathbb{Z},
}
\]

or

\[
\boxed{
\{0,1\}.
}
\]

Discrete restrictions lead to integer and combinatorial optimization.

%%%%%%%%%%%%%%%%%%%%%%%%%%%%%%%%%%%%%%%%%%%%%%%%%%%%%%%%%%%%
\subsection{Operations Research Perspective}
\label{subsec:operations-research-perspective}
%%%%%%%%%%%%%%%%%%%%%%%%%%%%%%%%%%%%%%%%%%%%%%%%%%%%%%%%%%%%

Operations research uses mathematical models to support decisions about
systems involving limited resources.

Typical questions include:

\[
\boxed{
\text{What should be produced?}
}
\]

\[
\boxed{
\text{How should resources be allocated?}
}
\]

\[
\boxed{
\text{Which route should be chosen?}
}
\]

\[
\boxed{
\text{When should tasks be scheduled?}
}
\]

and

\[
\boxed{
\text{How should uncertainty be managed?}
}
\]

Optimization provides the mathematical core of many such models.

%%%%%%%%%%%%%%%%%%%%%%%%%%%%%%%%%%%%%%%%%%%%%%%%%%%%%%%%%%%%
\subsection{Optimization Modeling Workflow}
\label{subsec:optimization-modeling-workflow}
%%%%%%%%%%%%%%%%%%%%%%%%%%%%%%%%%%%%%%%%%%%%%%%%%%%%%%%%%%%%

A practical optimization workflow is:

\begin{enumerate}
    \item define the decision to be made;
    \item choose decision variables;
    \item identify the objective;
    \item determine all physical, logical, financial, or mathematical
          constraints;
    \item define parameters and units;
    \item write the mathematical model;
    \item check dimensional consistency;
    \item determine whether the problem is linear, nonlinear, convex, or
          discrete;
    \item verify that the feasible set is meaningful;
    \item choose an appropriate solution method;
    \item check feasibility of the computed solution;
    \item check optimality conditions when possible;
    \item perform sensitivity analysis;
    \item interpret the solution in the original application.
\end{enumerate}

%%%%%%%%%%%%%%%%%%%%%%%%%%%%%%%%%%%%%%%%%%%%%%%%%%%%%%%%%%%%
\subsection{Choosing an Optimization Method}
\label{subsec:choosing-optimization-method}
%%%%%%%%%%%%%%%%%%%%%%%%%%%%%%%%%%%%%%%%%%%%%%%%%%%%%%%%%%%%

For a smooth unconstrained problem, consider:

\[
\boxed{
\text{gradient or Newton-type methods}.
}
\]

For equality constraints, consider:

\[
\boxed{
\text{Lagrange multipliers}.
}
\]

For smooth inequality constraints, consider:

\[
\boxed{
\text{KKT-based methods}.
}
\]

For linear objectives and linear constraints, use:

\[
\boxed{
\text{linear programming}.
}
\]

For integer decisions, use:

\[
\boxed{
\text{integer programming}.
}
\]

For convex nonlinear problems, use:

\[
\boxed{
\text{convex optimization methods}.
}
\]

%%%%%%%%%%%%%%%%%%%%%%%%%%%%%%%%%%%%%%%%%%%%%%%%%%%%%%%%%%%%
\subsection{Common Errors}
\label{subsec:optimization-fundamentals-common-errors}
%%%%%%%%%%%%%%%%%%%%%%%%%%%%%%%%%%%%%%%%%%%%%%%%%%%%%%%%%%%%

\subsubsection{Optimizing Before Defining the Variables}

Every model should state clearly what each decision variable represents.

Otherwise the mathematical solution may have no interpretable meaning.

%%%%%%%%%%%%%%%%%%%%%%%%%%%%%%%%%%%%%%%%%%%%%%%%%%%%%%%%%%%%
\subsubsection{Forgetting Units}

An objective such as

\[
3x+5y
\]

should correspond to meaningful and compatible quantities.

Unit consistency helps detect incorrect models.

%%%%%%%%%%%%%%%%%%%%%%%%%%%%%%%%%%%%%%%%%%%%%%%%%%%%%%%%%%%%
\subsubsection{Confusing Parameters with Decision Variables}

Parameters are treated as given data.

Decision variables are chosen by the optimization.

%%%%%%%%%%%%%%%%%%%%%%%%%%%%%%%%%%%%%%%%%%%%%%%%%%%%%%%%%%%%
\subsubsection{Ignoring Feasibility}

A point with an excellent objective value is irrelevant if it violates a
constraint.

Optimization always compares feasible decisions.

%%%%%%%%%%%%%%%%%%%%%%%%%%%%%%%%%%%%%%%%%%%%%%%%%%%%%%%%%%%%
\subsubsection{Assuming Every Stationary Point Is a Minimum}

The condition

\[
\nabla f=0
\]

is only a first-order necessary condition for an unconstrained interior
minimum.

The point may be a maximum or saddle.

%%%%%%%%%%%%%%%%%%%%%%%%%%%%%%%%%%%%%%%%%%%%%%%%%%%%%%%%%%%%
\subsubsection{Assuming Positive Semidefinite Hessian Always Gives a Strict Minimum}

Positive semidefiniteness can be insufficient for strict local
optimality.

Higher-order behavior may matter in degenerate cases.

%%%%%%%%%%%%%%%%%%%%%%%%%%%%%%%%%%%%%%%%%%%%%%%%%%%%%%%%%%%%
\subsubsection{Ignoring Boundaries}

A constrained optimum can occur at a boundary point where

\[
\nabla f\neq0.
\]

The feasible geometry must be considered.

%%%%%%%%%%%%%%%%%%%%%%%%%%%%%%%%%%%%%%%%%%%%%%%%%%%%%%%%%%%%
\subsubsection{Using Lagrange Multipliers Without Constraint Regularity}

Multiplier equations can fail to characterize an optimum if the
constraint gradients are degenerate.

Constraint qualifications matter.

%%%%%%%%%%%%%%%%%%%%%%%%%%%%%%%%%%%%%%%%%%%%%%%%%%%%%%%%%%%%
\subsubsection{Forgetting Dual Feasibility}

Under the convention

\[
g_i(x)\leq0,
\]

KKT inequality multipliers must satisfy

\[
\boxed{
\mu_i\geq0.
}
\]

Changing the sign convention changes the multiplier sign convention.

%%%%%%%%%%%%%%%%%%%%%%%%%%%%%%%%%%%%%%%%%%%%%%%%%%%%%%%%%%%%
\subsubsection{Forgetting Complementary Slackness}

For each inequality,

\[
\boxed{
\mu_i g_i(x^*)=0.
}
\]

Inactive constraints have zero multiplier.

%%%%%%%%%%%%%%%%%%%%%%%%%%%%%%%%%%%%%%%%%%%%%%%%%%%%%%%%%%%%
\subsubsection{Assuming KKT Conditions Always Prove Global Optimality}

For general nonconvex problems, KKT points may be local minima, maxima,
or saddle-type candidates.

Convexity is what often turns KKT conditions into sufficient conditions.

%%%%%%%%%%%%%%%%%%%%%%%%%%%%%%%%%%%%%%%%%%%%%%%%%%%%%%%%%%%%
\subsubsection{Confusing Local and Global Optimality}

A numerical solver may return a stationary or locally optimal point.

This does not automatically prove global optimality for a nonconvex
problem.

%%%%%%%%%%%%%%%%%%%%%%%%%%%%%%%%%%%%%%%%%%%%%%%%%%%%%%%%%%%%
\subsubsection{Using Matrix Inverses Numerically}

Expressions such as

\[
Q^{-1}c
\]

are mathematically convenient.

Numerically, solve

\[
Qx=c
\]

using an appropriate linear-system method instead of explicitly computing
the inverse.

%%%%%%%%%%%%%%%%%%%%%%%%%%%%%%%%%%%%%%%%%%%%%%%%%%%%%%%%%%%%
\subsubsection{Ignoring Scaling and Conditioning}

A theoretically correct method can behave poorly on badly scaled
problems.

Scaling should be treated as part of the numerical model.

%%%%%%%%%%%%%%%%%%%%%%%%%%%%%%%%%%%%%%%%%%%%%%%%%%%%%%%%%%%%
\subsubsection{Stopping Only Because Iterates Barely Move}

Small step size does not necessarily imply near-optimality.

A solver may stagnate.

Gradient, feasibility, and optimality residuals should also be examined.

%%%%%%%%%%%%%%%%%%%%%%%%%%%%%%%%%%%%%%%%%%%%%%%%%%%%%%%%%%%%
\subsection{Applications}
\label{subsec:optimization-fundamentals-applications}
%%%%%%%%%%%%%%%%%%%%%%%%%%%%%%%%%%%%%%%%%%%%%%%%%%%%%%%%%%%%

\begin{applicationbox}

\textbf{Operations Research.}
Optimization determines production plans, schedules, routes, inventory
policies, and resource allocations.

\medskip

\textbf{Engineering.}
Design variables are selected to minimize weight, cost, energy, or error
while satisfying physical constraints.

\medskip

\textbf{Machine Learning.}
Model parameters are estimated by minimizing loss functions, often with
regularization.

\medskip

\textbf{Statistics.}
Maximum likelihood, least squares, penalized estimation, and Bayesian MAP
estimation are optimization problems.

\medskip

\textbf{Economics.}
Consumers, firms, and planners optimize utility, profit, cost, welfare,
and allocation subject to constraints.

\medskip

\textbf{Finance.}
Portfolio optimization balances expected return, risk, and investment
constraints.

\medskip

\textbf{Control Theory.}
Optimal control selects inputs or trajectories that minimize a performance
criterion over time.

\medskip

\textbf{Science.}
Parameter estimation and inverse problems are often formulated as
optimization models.

\end{applicationbox}

%%%%%%%%%%%%%%%%%%%%%%%%%%%%%%%%%%%%%%%%%%%%%%%%%%%%%%%%%%%%
\subsection{Exercises}
\label{subsec:optimization-fundamentals-exercises}
%%%%%%%%%%%%%%%%%%%%%%%%%%%%%%%%%%%%%%%%%%%%%%%%%%%%%%%%%%%%

\begin{exercise}[1]
Define an optimization problem.
\end{exercise}

\begin{exercise}[1]
Define a decision variable.
\end{exercise}

\begin{exercise}[1]
Define an objective function.
\end{exercise}

\begin{exercise}[1]
Define the feasible set.
\end{exercise}

\begin{exercise}[1]
Define a global minimizer.
\end{exercise}

\begin{exercise}[1]
Define a local minimizer.
\end{exercise}

\begin{exercise}[1]
Define a convex set.
\end{exercise}

\begin{exercise}[1]
Define a convex function.
\end{exercise}

\begin{exercise}[1]
Define an active inequality constraint.
\end{exercise}

\begin{exercise}[1]
State the four basic KKT conditions.
\end{exercise}

\begin{exercise}[2]
Find the stationary point of

\[
f(x)=x^2-8x+20.
\]
\end{exercise}

\begin{exercise}[2]
Classify the stationary point in the previous problem.
\end{exercise}

\begin{exercise}[2]
Find all stationary points of

\[
f(x,y)
=
x^2+y^2-2x+6y.
\]
\end{exercise}

\begin{exercise}[2]
Compute the Hessian of

\[
f(x,y)
=
3x^2+2xy+4y^2.
\]
\end{exercise}

\begin{exercise}[2]
Determine whether

\[
Q
=
\begin{pmatrix}
2&0\\
0&5
\end{pmatrix}
\]

is positive definite.
\end{exercise}

\begin{exercise}[2]
Determine whether the set

\[
\{
(x,y):
x+y\leq1
\}
\]

is convex.
\end{exercise}

\begin{exercise}[2]
Determine whether

\[
f(x)=x^2
\]

is convex.
\end{exercise}

\begin{exercise}[2]
Determine whether

\[
f(x)=-x^2
\]

is convex.
\end{exercise}

\begin{exercise}[2]
Identify which constraints are active at

\[
(x,y)=(1,2)
\]

for

\[
x\leq1,
\qquad
y\leq3,
\qquad
x+y\leq3.
\]
\end{exercise}

\begin{exercise}[3]
Minimize

\[
f(x,y)=x^2+y^2
\]

subject to

\[
x+y=4
\]

using Lagrange multipliers.
\end{exercise}

\begin{exercise}[3]
Minimize

\[
f(x,y)=x^2+4y^2
\]

subject to

\[
x+2y=6.
\]
\end{exercise}

\begin{exercise}[3]
Find the KKT point for

\[
\min_x
(x-5)^2
\]

subject to

\[
x\leq2.
\]
\end{exercise}

\begin{exercise}[3]
Find the KKT point for

\[
\min_x
x^2
\]

subject to

\[
x\geq3.
\]
\end{exercise}

\begin{exercise}[3]
Verify complementary slackness in the previous problem.
\end{exercise}

\begin{exercise}[3]
Derive the normal equations for least squares.
\end{exercise}

\begin{exercise}[3]
Show that

\[
f(x)
=
\frac12x^TQx+c^Tx
\]

has gradient

\[
Qx+c
\]

when \(Q\) is symmetric.
\end{exercise}

\begin{exercise}[3]
Show that if \(Q\succ0\), then the quadratic objective has a unique
global minimizer.
\end{exercise}

\begin{exercise}[3]
Prove that every local minimum of a convex function on a convex set is a
global minimum.
\end{exercise}

\begin{exercise}[3]
Prove that a strictly convex function has at most one minimizer.
\end{exercise}

\begin{exercise}[3]
Take one gradient-descent step for

\[
f(x,y)
=
x^2+2y^2
\]

from

\[
(x_0,y_0)=(2,1)
\]

using step size

\[
\alpha=0.1.
\]
\end{exercise}

\begin{exercise}[3]
Derive the Newton step from the quadratic Taylor approximation.
\end{exercise}

\begin{exercise}[3]
Apply Newton's method to

\[
f(x)=x^4-4x^2+5
\]

from a specified initial point.
\end{exercise}

\begin{exercise}[3]
Explain why poor scaling can slow gradient descent.
\end{exercise}

\begin{exercise}[3]
Formulate a production problem with two products, three resource
constraints, and nonnegative decision variables.
\end{exercise}

\begin{exercise}[3]
Formulate a portfolio model with weights satisfying

\[
\sum_iw_i=1.
\]
\end{exercise}

\begin{exercise}[3]
Formulate a least-squares regression problem as an optimization problem.
\end{exercise}

\begin{exercise}[3]
Explain how \(\ell_2\) regularization changes a least-squares objective.
\end{exercise}

\begin{exercise}[3]
Explain why \(\ell_1\) regularization can produce sparse solutions.
\end{exercise}

\begin{exercise}[4]
Prove the first-order characterization of differentiable convex
functions.
\end{exercise}

\begin{exercise}[4]
Prove the Hessian criterion for twice-differentiable convex functions.
\end{exercise}

\begin{exercise}[4]
Study strong convexity and derive its basic inequalities.
\end{exercise}

\begin{exercise}[4]
Study Lipschitz continuity of gradients.
\end{exercise}

\begin{exercise}[4]
Derive convergence bounds for gradient descent on smooth strongly convex
functions.
\end{exercise}

\begin{exercise}[4]
Study Newton's local quadratic convergence.
\end{exercise}

\begin{exercise}[4]
Derive Wolfe or Armijo line-search conditions.
\end{exercise}

\begin{exercise}[4]
Study trust-region optimality conditions.
\end{exercise}

\begin{exercise}[4]
Derive equality-constrained second-order optimality conditions.
\end{exercise}

\begin{exercise}[4]
Study LICQ and MFCQ constraint qualifications.
\end{exercise}

\begin{exercise}[4]
Study Slater's condition for convex problems.
\end{exercise}

\begin{exercise}[4]
Derive the KKT conditions from tangent-cone arguments.
\end{exercise}

\begin{exercise}[4]
Study normal cones and first-order constrained optimality.
\end{exercise}

\begin{exercise}[4]
Investigate exact penalty functions.
\end{exercise}

\begin{exercise}[4]
Study logarithmic barrier methods.
\end{exercise}

\begin{exercise}[4]
Study augmented Lagrangian methods.
\end{exercise}

\begin{exercise}[4]
Investigate sensitivity of the optimal value to constraint perturbations.
\end{exercise}

\begin{exercise}[4]
Study regularization as constrained optimization through norm balls.
\end{exercise}

\begin{exercise}[4]
Compare deterministic, stochastic, and robust optimization.
\end{exercise}

%%%%%%%%%%%%%%%%%%%%%%%%%%%%%%%%%%%%%%%%%%%%%%%%%%%%%%%%%%%%
\subsection{Conceptual Summary}
\label{subsec:optimization-fundamentals-summary}
%%%%%%%%%%%%%%%%%%%%%%%%%%%%%%%%%%%%%%%%%%%%%%%%%%%%%%%%%%%%

A general optimization problem is

\[
\boxed{
\begin{aligned}
\min_{\mathbf{x}}
\quad&
f(\mathbf{x})
\\
\text{subject to}
\quad&
g_i(\mathbf{x})\leq0,
\\
&
h_j(\mathbf{x})=0.
\end{aligned}
}
\]

The feasible set contains all points satisfying the constraints.

For an unconstrained smooth local minimizer,

\[
\boxed{
\nabla f(\mathbf{x}^*)=0.
}
\]

If additionally

\[
\boxed{
\nabla^2f(\mathbf{x}^*)\succ0,
}
\]

the point is a strict local minimizer.

For convex problems,

\[
\boxed{
\text{local minimum}
=
\text{global minimum}.
}
\]

Equality constraints lead to the Lagrangian

\[
\boxed{
\mathcal{L}
=
f
+
\sum_j
\lambda_jh_j.
}
\]

Inequality constraints lead to the KKT conditions:

\[
\boxed{
\text{stationarity},
}
\]

\[
\boxed{
\text{primal feasibility},
}
\]

\[
\boxed{
\text{dual feasibility},
}
\]

and

\[
\boxed{
\text{complementary slackness}.
}
\]

Numerical optimization methods generate iterates such as

\[
\boxed{
x_{k+1}
=
x_k
-
\alpha_k
\nabla f(x_k).
}
\]

Optimization therefore combines

\[
\boxed{
\text{calculus}
+
\text{linear algebra}
+
\text{geometry}
+
\text{modeling}
+
\text{numerical algorithms}.
}
\]

%%%%%%%%%%%%%%%%%%%%%%%%%%%%%%%%%%%%%%%%%%%%%%%%%%%%%%%%%%%%
\subsection{Section Connections}
\label{subsec:optimization-fundamentals-connections}
%%%%%%%%%%%%%%%%%%%%%%%%%%%%%%%%%%%%%%%%%%%%%%%%%%%%%%%%%%%%

\chapterconnection{
Optimization Fundamentals establishes the language used throughout
Chapter~10: decision variables, objectives, feasible sets, convexity,
Lagrange multipliers, KKT conditions, and numerical search methods.
These ideas connect directly with the constrained systems encountered in
Chapter~9 and with estimation problems in Statistics. The next section,
Linear Programming, specializes optimization to linear objectives and
linear constraints, leading to polyhedral geometry, basic feasible
solutions, the simplex method, duality, sensitivity analysis, and some of
the most important models in operations research.
}

%%%%%%%%%%%%%%%%%%%%%%%%%%%%%%%%%%%%%%%%%%%%%%%%%%%%%%%%%%%%
% SECTION SOURCE TRACKING
%%%%%%%%%%%%%%%%%%%%%%%%%%%%%%%%%%%%%%%%%%%%%%%%%%%%%%%%%%%%

%------------------------------------------------------------
% 10.1 Optimization Fundamentals
%------------------------------------------------------------
%
% Source basis:
%   - Standard mathematical optimization theory.
%   - Standard nonlinear programming theory.
%   - Standard introductory operations-research modeling.
%
% Used for:
%   - optimization terminology;
%   - decision variables;
%   - objectives and constraints;
%   - feasible sets;
%   - local and global optima;
%   - unconstrained and constrained optimization;
%   - optimization modeling;
%   - gradients and directional derivatives;
%   - first-order necessary conditions;
%   - Hessians;
%   - second-order conditions;
%   - positive definiteness;
%   - quadratic optimization;
%   - least squares;
%   - existence and coercivity;
%   - convex sets and convex functions;
%   - strict convexity;
%   - first- and second-order convexity tests;
%   - equality-constrained optimization;
%   - Lagrange multipliers;
%   - inequality constraints;
%   - active constraints;
%   - KKT conditions;
%   - complementary slackness;
%   - constraint qualifications;
%   - multiplier sensitivity;
%   - regularization;
%   - penalty and barrier methods;
%   - gradient descent;
%   - Newton's method;
%   - line-search and trust-region previews;
%   - quasi-Newton methods;
%   - stopping criteria;
%   - scaling and conditioning;
%   - deterministic, stochastic, multiobjective, and discrete
%     optimization previews;
%   - operations-research modeling;
%   - applications and exercises.
%
% Notes:
%   - Linear Programming is developed next in Section 10.2.
%   - Nonlinear Programming and Convex Optimization develop the smooth
%     and convex theory more deeply later in Chapter 10.
%   - Numerical Optimization is developed more deeply in Chapter 11.
%
%%%%%%%%%%%%%%%%%%%%%%%%%%%%%%%%%%%%%%%%%%%%%%%%%%%%%%%%%%%%